\documentclass[12pt,a4paper,oneside]{book}

% =====================================================================
% PACOTES ESSENCIAIS
% =====================================================================
\usepackage[portuguese]{babel}
\usepackage[T1]{fontenc}
\usepackage{amsmath}
\usepackage{amssymb}
\usepackage{amsthm}
\usepackage{geometry}
\usepackage{fancyhdr}
\usepackage{graphicx}
\usepackage{xcolor}
\usepackage{listings}
\usepackage{hyperref}
\usepackage{booktabs}
\usepackage{array}
\usepackage{multirow}
\usepackage{makeidx}
\usepackage{tikz}
\usepackage{mathtools}

% =====================================================================
% CONFIGURAÇÕES DE GEOMETRIA
% =====================================================================
\geometry{
    left=2.5cm,
    right=2.5cm,
    top=3cm,
    bottom=3cm,
    headheight=14pt
}

% =====================================================================
% CONFIGURAÇÕES DE CABEÇALHO E RODAPÉ
% =====================================================================
\pagestyle{fancy}
\fancyhf{}
\fancyhead[L]{\textit{Capítulo 2: Análise Dimensional}}
\fancyhead[R]{\thepage}
\fancyfoot[C]{\small Princípios de Modelagem Matemática}
\renewcommand{\headrulewidth}{0.4pt}
\renewcommand{\footrulewidth}{0.4pt}

% =====================================================================
% DEFINIÇÃO DE TEOREMAS E AMBIENTES
% =====================================================================
\theoremstyle{definition}
\newtheorem{definition}{Definição}[chapter]
\newtheorem{theorem}[definition]{Teorema}
\newtheorem{lemma}[definition]{Lema}
\newtheorem{corollary}[definition]{Corolário}
\newtheorem{example}[definition]{Exemplo}
\newtheorem{remark}[definition]{Observação}
\newtheorem{problem}{Problema}

% =====================================================================
% CONFIGURAÇÕES DE CÓDIGO
% =====================================================================
\lstset{
    language=Python,
    basicstyle=\ttfamily\small,
    keywordstyle=\color{blue},
    commentstyle=\color{gray},
    stringstyle=\color{red},
    breaklines=true,
    showstringspaces=false,
    numbers=left,
    numberstyle=\tiny,
    numbersep=5pt,
    frame=single,
    frameround=tttt,
    rulecolor=\color{black}
}

% =====================================================================
% MACROS PERSONALIZADAS
% =====================================================================
\newcommand{\dimension}[1]{$[#1]$}
\newcommand{\dimhom}{\text{dimensionalmente homogêneo}}
\newcommand{\dimless}{\text{adimensional}}
\newcommand{\pithm}{\text{Teorema de Buckingham Pi}}
\newcommand{\todo}[1]{\textcolor{red}{\textbf{TODO: #1}}}

% =====================================================================
% CORES
% =====================================================================
\definecolor{darkblue}{RGB}{0,0,139}
\definecolor{lightblue}{RGB}{173,216,230}
\definecolor{darkgreen}{RGB}{0,100,0}
\definecolor{lightgreen}{RGB}{144,238,144}

% =====================================================================
% DOCUMENTO
% =====================================================================
\title{\textbf{Capítulo 2: Análise Dimensional} \\[0.5cm]
        \large Princípios de Modelagem Matemática}
\author{Baseado em Dym, C. L. (2004) \\
        \small Principles of Mathematical Modeling, 2nd Edition}
\date{2026-04-14}

\makeindex

\begin{document}

% =====================================================================
% FOLHA DE ROSTO
% =====================================================================
\maketitle

% =====================================================================
% SUMÁRIO
% =====================================================================
\tableofcontents
\newpage

% =====================================================================
% CAPÍTULO 2: ANÁLISE DIMENSIONAL
% =====================================================================
\chapter{Análise Dimensional}

\section{Introdução}

A \textbf{Análise Dimensional} é uma ferramenta fundamental no processo de modelagem matemática que garante a consistência lógica e física das equações. Este capítulo apresenta metodicamente os conceitos e técnicas para:

\begin{itemize}
    \item Validar modelos matemáticos antes de realizar cálculos custosos
    \item Reduzir drasticamente o número de variáveis necessárias em experimentos
    \item Entender as relações fundamentais entre grandezas físicas
    \item Detectar erros na formulação de problemas
\end{itemize}

\subsection{Motivação: Economia de Experimentos}

Um exemplo motivador: suponha que desejamos modelar a força de arrasto em um misturador:

\begin{equation}
    F_D = F_D(V, d, \rho, \mu)
\end{equation}

onde:
\begin{align}
    V &= \text{velocidade da lâmina} \\
    d &= \text{largura da lâmina} \\
    \rho &= \text{densidade do fluido} \\
    \mu &= \text{viscosidade do fluido}
\end{align}

Se realizássemos experimentos testando 3 valores para cada variável:
\begin{equation}
    \text{Número de experimentos} = 3^5 = 243
\end{equation}

Com análise dimensional, reduzimos para apenas 2 grupos adimensionais, necessitando:
\begin{equation}
    \text{Número de experimentos} = 3^2 = 9 \quad \text{(redução de } 96\%\text{!)}
\end{equation}

% =====================================================================
% SEÇÃO 1: DIMENSÕES E UNIDADES
% =====================================================================
\section{Dimensões e Unidades}

\subsection{Definiçōes Fundamentais}

\begin{definition}[Dimensão]
A \textbf{dimensão} de uma quantidade física é sua natureza fundamental expressa em termos de grandezas primitivas, frequentemente denotadas como:
\begin{equation}
    \text{Dimensão} = M^a L^b T^c \Theta^d \cdots
\end{equation}
onde $M$ (massa), $L$ (comprimento), $T$ (tempo), $\Theta$ (temperatura) são as sete dimensões fundamentais do Sistema Internacional.
\end{definition}

\begin{definition}[Unidade]
A \textbf{unidade} é a expressão numérica de uma dimensão em relação a um padrão estabelecido. A unidade depende da escolha do sistema de medição, mas a dimensão é universal.
\end{definition}

\begin{example}[Dimensão vs. Unidade]
Para a velocidade:
\begin{align}
    \text{Dimensão} &: [v] = LT^{-1} \\
    \text{Unidades possíveis} &: \text{m/s}, \text{ ft/s}, \text{ mi/hr}, \text{ km/h}
\end{align}
\end{example}

\subsection{Dimensões de Quantidades Comuns}

A Tabela \ref{tab:dimensions} resume as dimensões de quantidades físicas comuns.

\begin{table}[h]
\centering
\caption{Dimensões de Quantidades Físicas Comuns}
\label{tab:dimensions}
\begin{tabular}{ll}
\toprule
\textbf{Quantidade} & \textbf{Dimensão} \\
\midrule
Comprimento & $L$ \\
Tempo & $T$ \\
Massa & $M$ \\
Velocidade & $LT^{-1}$ \\
Aceleração & $LT^{-2}$ \\
Força & $MLT^{-2}$ \\
Energia/Trabalho & $ML^2T^{-2}$ \\
Potência & $ML^2T^{-3}$ \\
Pressão & $ML^{-1}T^{-2}$ \\
Viscosidade Dinâmica & $ML^{-1}T^{-1}$ \\
Densidade & $ML^{-3}$ \\
Frequência & $T^{-1}$ \\
\bottomrule
\end{tabular}
\end{table}

% =====================================================================
% SEÇÃO 2: HOMOGENEIDADE DIMENSIONAL
% =====================================================================
\section{Homogeneidade Dimensional}

\subsection{O Princípio Fundamental}

\begin{theorem}[Homogeneidade Dimensional]
Em qualquer equação que descreva um fenômeno físico, todos os termos independentes devem ter as mesmas dimensões. Uma equação que satisfaz este requisito é chamada \textbf{dimensionalmente homogênea} ou \textbf{racional}.
\end{theorem}

Este princípio é tão fundamental que frequentemente utilizamos a notação de parênteses quadrados $[~\cdot~]$ para indicar ``a dimensão de'':

\begin{equation}
    [A] = [B] = [C] \quad \text{para a equação } A = B + C
\end{equation}

\subsection{Exemplos de Homogeneidade}

\subsubsection{Exemplo 1: Período do Pêndulo}

A fórmula clássica para o período de um pêndulo simples é:
\begin{equation}
    T_0 = 2\pi\sqrt{\frac{l}{g}}
\end{equation}

Verificamos a homogeneidade:
\begin{align}
    [T_0] &= T \\
    \left[\sqrt{\frac{l}{g}}\right] &= \sqrt{\frac{L}{LT^{-2}}} = \sqrt{T^2} = T \quad \checkmark
\end{align}

\subsubsection{Exemplo 2: Equação Inválida}

Considere a equação:
\begin{equation}
    T_0 = 2\pi\sqrt{\frac{l}{g}} + 5 \text{ segundos}
\end{equation}

Esta equação é \textbf{dimensionalmente inconsistente} porque:
\begin{align}
    \text{Termo 1} &: [2\pi\sqrt{l/g}] = T \\
    \text{Termo 2} &: [5 \text{ s}] = T \\
    \text{Mas } &: T + T \neq \text{(adicionar um vetor constante a uma função)}
\end{align}

A equação viola o princípio de homogeneidade dimensional e deve ser rejeitada.

% =====================================================================
% SEÇÃO 3: POR QUE ANÁLISE DIMENSIONAL?
% =====================================================================
\section{Por Que Fazer Análise Dimensional?}

\subsection{Benefícios Práticos}

A análise dimensional oferece vários benefícios críticos para o modelador:

\begin{enumerate}
    \item \textbf{Detecção de Erros}: Erros algébricos frequentemente violam a homogeneidade dimensional
    \item \textbf{Economia de Experimentos}: Redução exponencial no número de testes necessários
    \item \textbf{Compreensão Física}: Revela quais variáveis são realmente importantes
    \item \textbf{Generalização}: Resultados em uma unidade aplicam-se a outras via grupos adimensionais
    \item \textbf{Validação}: Compara modelos teóricos com observações experimentais
\end{enumerate}

\subsection{Caso de Estudo: Misturador de Manteiga de Amendoim}

Considere o problema de projetar um misturador industrial para manteiga de amendoim. Queremos modelar como a força de arrasto na lâmina ($F_D$) depende de:

\begin{itemize}
    \item Velocidade da lâmina: $V$ [dimensão: $LT^{-1}$]
    \item Largura da lâmina: $d$ [dimensão: $L$]
    \item Densidade da manteiga: $\rho$ [dimensão: $ML^{-3}$]
    \item Viscosidade da manteiga: $\mu$ [dimensão: $ML^{-1}T^{-1}$]
\end{itemize}

\textbf{Abordagem ingênua} (sem análise dimensional):
\begin{align}
    \text{Variáveis} &= 5 \quad (F_D, V, d, \rho, \mu) \\
    \text{Valores por variável} &= 3 \\
    \text{Experimentos necessários} &= 3^5 = 243
\end{align}

\textbf{Abordagem inteligente} (com análise dimensional):
\begin{align}
    \text{Variáveis} &= 5 \\
    \text{Dimensões fundamentais} &= 3 \quad (M, L, T) \\
    \text{Grupos adimensionais} &= 5 - 3 = 2 \\
    \text{Experimentos necessários} &\approx 3^2 = 9
\end{align}

\textbf{Economia}: $243 - 9 = 234$ experimentos economizados ($\approx 96\%$ de redução)!

% =====================================================================
% SEÇÃO 4: COMO FAZER ANÁLISE DIMENSIONAL
% =====================================================================
\section{Como Fazer Análise Dimensional?}

Existem dois métodos principais para aplicar análise dimensional. O primeiro é informal e intuitivo; o segundo é formal e sistemático.

\subsection{Método 1: Método Básico}

O método básico de análise dimensional é aplicado através dos seguintes passos:

\begin{enumerate}
    \item \textbf{Listar variáveis}: Liste todas as variáveis e parâmetros do problema com suas dimensões
    \item \textbf{Antecipar efeitos}: Para cada variável, determine se um aumento causa aumento ou diminuição na variável dependente
    \item \textbf{Identificar dependente}: Escolha uma variável como dependente das demais
    \item \textbf{Formular funcionalmente}: Expresse a dependência como uma equação funcional
    \item \textbf{Eliminar dimensões}: Escolha e elimine uma dimensão primária de cada vez
    \item \textbf{Repetir}: Repita o processo até obter uma equação adimensional
    \item \textbf{Validar}: Revise se o comportamento esperado é fisicamente razoável
\end{enumerate}

\subsubsection{Exemplo: Velocidade de Corpo em Queda}

\textbf{Problema}: Um corpo cai de uma altura $h$ sujeito à aceleração da gravidade $g$. Encontre sua velocidade $V$ ao chegar ao solo usando análise dimensional.

\textbf{Solução - Passo a Passo}:

\begin{enumerate}
    \item \textbf{Variáveis e Dimensões}:
    \begin{align}
        V &: \text{velocidade} \quad [V] = LT^{-1} \\
        g &: \text{aceleração} \quad [g] = LT^{-2} \\
        h &: \text{altura} \quad [h] = L \\
        m &: \text{massa} \quad [m] = M
    \end{align}

    \item \textbf{Função}:
    \begin{equation}
        V = V(g, h, m)
    \end{equation}

    \item \textbf{Forma Paramétrica}:
    \begin{equation}
        V = C \cdot g^a \cdot h^b \cdot m^c
    \end{equation}
    onde $C$ é uma constante adimensional.

    \item \textbf{Análise Dimensional}:
    \begin{equation}
        [LT^{-1}] = [LT^{-2}]^a \cdot [L]^b \cdot [M]^c
    \end{equation}

    \begin{equation}
        LT^{-1} = L^a T^{-2a} \cdot L^b \cdot M^c = M^c L^{a+b} T^{-2a}
    \end{equation}

    \item \textbf{Igualando Expoentes}:
    \begin{align}
        M: & \quad 0 = c \quad \Rightarrow \quad c = 0 \quad \text{(massa não influencia!)} \\
        L: & \quad 1 = a + b \\
        T: & \quad -1 = -2a \quad \Rightarrow \quad a = \frac{1}{2}
    \end{align}

    \item \textbf{Resolvendo}:
    \begin{equation}
        b = 1 - a = 1 - \frac{1}{2} = \frac{1}{2}
    \end{equation}

    \item \textbf{Resultado}:
    \begin{equation}
        V = C \sqrt{gh}
    \end{equation}

    Da física conhecemos $C = \sqrt{2}$, portanto:
    \begin{equation}
        V = \sqrt{2gh}
    \end{equation}

    \item \textbf{Validação}: A equação faz sentido fisicamente:
    \begin{itemize}
        \item Maior $g$ $\Rightarrow$ maior $V$ \checkmark
        \item Maior $h$ $\Rightarrow$ maior $V$ \checkmark
        \item Massa não afeta $V$ \checkmark (confirmado experimentalmente por Galileu)
    \end{itemize}
\end{enumerate}

\subsection{Método 2: Teorema de Buckingham Pi}

O Teorema de Buckingham Pi é uma abordagem formal e sistemática para análise dimensional.

\begin{theorem}[Buckingham Pi]
Uma equação dimensionalmente homogênea envolvendo $n$ variáveis e $m$ dimensões primárias independentes pode ser reduzida a uma relação entre $n - m$ produtos adimensionais (chamados ``grupos Pi'').

Se as variáveis são $A_1, A_2, \ldots, A_n$, então:

\begin{equation}
    \Pi_1 = \Phi(\Pi_2, \Pi_3, \ldots, \Pi_{n-m})
\end{equation}

ou equivalentemente:

\begin{equation}
    \Phi(\Pi_1, \Pi_2, \ldots, \Pi_{n-m}) = 0
\end{equation}

onde cada $\Pi_i$ é um grupo adimensional.
\end{theorem}

\subsubsection{Aplicação do Teorema Pi}

O processo de aplicação é sistemático:

\begin{enumerate}
    \item Identificar as $n$ variáveis do problema
    \item Contar as $m$ dimensões fundamentais
    \item Calcular número de grupos: $n - m$
    \item Selecionar $m$ variáveis ``base'' que contenham todas as $m$ dimensões
    \item Construir grupos permutando as $n - m$ variáveis restantes com as base
    \item Resolver o sistema linear para encontrar os expoentes de cada grupo
\end{enumerate}

\subsubsection{Exemplo: Período do Pêndulo}

\textbf{Problema}: Encontrar os grupos adimensionais que relacionam o período de um pêndulo simples às suas variáveis características.

\textbf{Variáveis}:
\begin{align}
    l &: \text{comprimento} \quad [l] = L \\
    g &: \text{aceleração} \quad [g] = LT^{-2} \\
    m &: \text{massa} \quad [m] = M \\
    T_0 &: \text{período} \quad [T_0] = T \\
    \theta &: \text{amplitude} \quad [\theta] = \text{(adimensional)}
\end{align}

\textbf{Análise}:
\begin{align}
    n &= 5 \quad \text{(variáveis)} \\
    m &= 3 \quad \text{(dimensões: M, L, T)} \\
    n - m &= 2 \quad \text{(grupos adimensionais)}
\end{align}

\textbf{Grupos dimensionless}:

Escolhemos como base: $l$, $g$, $m$ (contêm M, L, T)

As variáveis restantes são: $T_0$, $\theta$

\textit{Grupo 1} (envolvendo $T_0$):
\begin{equation}
    \Pi_1 = T_0 \cdot l^{a_1} \cdot g^{b_1} \cdot m^{c_1}
\end{equation}

Para ser adimensional:
\begin{align}
    [T] \cdot L^{a_1} \cdot (LT^{-2})^{b_1} \cdot M^{c_1} &= 1 \\
    M^{c_1} \cdot L^{a_1 + b_1} \cdot T^{1-2b_1} &= 1
\end{align}

Igualando expoentes:
\begin{align}
    M: & \quad c_1 = 0 \\
    L: & \quad a_1 + b_1 = 0 \quad \Rightarrow \quad a_1 = -b_1 \\
    T: & \quad 1 - 2b_1 = 0 \quad \Rightarrow \quad b_1 = \frac{1}{2}
\end{align}

Portanto: $a_1 = -\frac{1}{2}$, $b_1 = \frac{1}{2}$, $c_1 = 0$

\begin{equation}
    \Pi_1 = T_0 \cdot l^{-1/2} \cdot g^{1/2} = \frac{T_0}{\sqrt{l/g}}
\end{equation}

\textit{Grupo 2} (envolvendo $\theta$):

Como $\theta$ já é adimensional:
\begin{equation}
    \Pi_2 = \theta
\end{equation}

\textbf{Relação Funcional}:
\begin{equation}
    \frac{T_0}{\sqrt{l/g}} = f(\theta)
\end{equation}

\textbf{Interpretação Física}:
\begin{equation}
    T_0 = \sqrt{\frac{l}{g}} \cdot f(\theta)
\end{equation}

Para pequenas amplitudes ($\theta \ll 1$), $f(\theta) \approx 2\pi$, resultando na conhecida fórmula:
\begin{equation}
    T_0 = 2\pi\sqrt{\frac{l}{g}}
\end{equation}

% =====================================================================
% SEÇÃO 5: SISTEMAS DE UNIDADES
% =====================================================================
\section{Sistemas de Unidades}

\subsection{Introdução}

Enquanto as dimensões são universais, as unidades dependem do sistema de medição adotado. É crítico ser consistente na escolha de unidades durante cálculos.

\subsection{Sistema SI vs. Sistema Britânico}

A Tabela \ref{tab:units} compara os dois sistemas mais comuns.

\begin{table}[h]
\centering
\caption{Comparação entre Sistemas SI e Britânico}
\label{tab:units}
\begin{tabular}{lcc}
\toprule
\textbf{Grandeza} & \textbf{SI} & \textbf{Britânico} \\
\midrule
Comprimento & metro (m) & pé (ft) \\
Massa & quilograma (kg) & slug \\
Tempo & segundo (s) & segundo (s) \\
Força & newton (N) & libra-força (lbf) \\
Pressão & pascal (Pa) & lbf/ft² \\
Potência & watt (W) & horsepower (hp) \\
\bottomrule
\end{tabular}
\end{table}

\subsection{Conversão de Unidades}

A conversão correta entre sistemas é crítica. Considere converter $65$ mi/hr para m/s:

\begin{equation}
65 \frac{\text{mi}}{\text{hr}} \times \frac{5280 \text{ ft}}{\text{mi}} \times \frac{0.3048 \text{ m}}{\text{ft}} \times \frac{1 \text{ hr}}{3600 \text{ s}}
\end{equation}

\begin{equation}
= 65 \times \frac{5280 \times 0.3048}{3600} \text{ m/s} = 29.06 \text{ m/s}
\end{equation}

\textbf{Verificação Dimensional}:
\begin{equation}
\frac{\text{mi}}{\text{hr}} \times \frac{\text{ft}}{\text{mi}} \times \frac{\text{m}}{\text{ft}} \times \frac{\text{hr}}{\text{s}} = \frac{\text{m}}{\text{s}} \quad \checkmark
\end{equation}

\subsection{Prefixos SI}

A Tabela \ref{tab:prefixes} lista os prefixos SI mais comuns.

\begin{table}[h]
\centering
\caption{Prefixos do Sistema Internacional}
\label{tab:prefixes}
\begin{tabular}{ccc}
\toprule
\textbf{Fator} & \textbf{Prefixo} & \textbf{Símbolo} \\
\midrule
$10^{12}$ & tera & T \\
$10^9$ & giga & G \\
$10^6$ & mega & M \\
$10^3$ & quilo & k \\
$10^0$ & (base) & \\
$10^{-3}$ & mili & m \\
$10^{-6}$ & micro & $\mu$ \\
$10^{-9}$ & nano & n \\
$10^{-12}$ & pico & p \\
\bottomrule
\end{tabular}
\end{table}

% =====================================================================
% SEÇÃO 6: EXEMPLOS RESOLVIDOS
% =====================================================================
\section{Exemplos Resolvidos}

\subsection{Exemplo 1: Velocidade do Som}

\begin{problem}
A velocidade do som $c$ em um gás depende da pressão $p$ e da densidade do gás $\rho$. Usando análise dimensional, determine a relação funcional.
\end{problem}

\textbf{Solução}:

\textbf{Variáveis e dimensões}:
\begin{align}
    [c] &= LT^{-1} \\
    [p] &= ML^{-1}T^{-2} \\
    [\rho] &= ML^{-3}
\end{align}

\textbf{Forma paramétrica}:
\begin{equation}
    c = C \cdot p^a \cdot \rho^b
\end{equation}

\textbf{Análise dimensional}:
\begin{equation}
    LT^{-1} = (ML^{-1}T^{-2})^a \cdot (ML^{-3})^b = M^{a+b} L^{-a-3b} T^{-2a}
\end{equation}

\textbf{Sistema de equações}:
\begin{align}
    M: & \quad 0 = a + b \quad \Rightarrow \quad b = -a \\
    L: & \quad 1 = -a - 3b = -a + 3a = 2a \quad \Rightarrow \quad a = \frac{1}{2} \\
    T: & \quad -1 = -2a \quad \Rightarrow \quad a = \frac{1}{2} \quad \checkmark
\end{align}

\textbf{Resultado}:
\begin{equation}
    c = C\sqrt{\frac{p}{\rho}}
\end{equation}

Fisicamente, $C$ é uma função do índice adiabático do gás, e a fórmula correta para gases ideais é:
\begin{equation}
    c = \sqrt{\gamma \frac{p}{\rho}}
\end{equation}

\subsection{Exemplo 2: Força de Arrasto em Fluido}

\begin{problem}
A força de arrasto $F_D$ sobre uma esfera movendo-se através de um fluido viscoso depende de:
- Velocidade: $V$
- Diâmetro: $d$
- Viscosidade: $\mu$

Encontre a relação funcional.
\end{problem}

\textbf{Solução}:

\textbf{Variáveis e dimensões}:
\begin{align}
    [F_D] &= MLT^{-2} \\
    [V] &= LT^{-1} \\
    [d] &= L \\
    [\mu] &= ML^{-1}T^{-1}
\end{align}

\textbf{Análise}:
\begin{align}
    n &= 4 \quad \text{(variáveis)} \\
    m &= 3 \quad \text{(dimensões)} \\
    n - m &= 1 \quad \text{(grupo adimensional)}
\end{align}

\textbf{Forma paramétrica}:
\begin{equation}
    F_D = C \cdot V^a \cdot d^b \cdot \mu^c
\end{equation}

\textbf{Análise dimensional}:
\begin{equation}
    MLT^{-2} = (LT^{-1})^a \cdot L^b \cdot (ML^{-1}T^{-1})^c = M^c L^{a+b-c} T^{-a-c}
\end{equation}

\textbf{Sistema de equações}:
\begin{align}
    M: & \quad 1 = c \\
    L: & \quad 1 = a + b - c = a + b - 1 \quad \Rightarrow \quad a + b = 2 \\
    T: & \quad -2 = -a - c = -a - 1 \quad \Rightarrow \quad a = 1
\end{align}

\textbf{Resolvendo}:
\begin{equation}
    b = 2 - a = 2 - 1 = 1
\end{equation}

\textbf{Resultado} (Lei de Stokes):
\begin{equation}
    F_D = 3\pi \mu V d
\end{equation}

% =====================================================================
% SEÇÃO 7: RESUMO E CHECKLIST
% =====================================================================
\section{Resumo e Checklist}

\subsection{Pontos-Chave}

\begin{enumerate}
    \item A \textbf{homogeneidade dimensional} é um requisito fundamental para qualquer modelo matemático válido
    \item \textbf{Dimensões} são universais; \textbf{unidades} dependem do sistema escolhido
    \item Análise dimensional reduz drasticamente o número de variáveis necessárias
    \item O \textbf{Método Básico} é intuitivo; o \textbf{Teorema de Buckingham Pi} é sistemático
    \item Sempre validar resultados com conhecimento físico
\end{enumerate}

\subsection{Checklist de Aplicação}

Ao aplicar análise dimensional, certifique-se de:

\begin{itemize}
    \item[$\square$] Listar todas as variáveis e parâmetros
    \item[$\square$] Determinar as dimensões de cada quantidade
    \item[$\square$] Contar o número de dimensões fundamentais
    \item[$\square$] Calcular o número de grupos dimensionless
    \item[$\square$] Escolher variáveis base apropriadas
    \item[$\square$] Formar grupos permutando as variáveis restantes
    \item[$\square$] Resolver o sistema linear para expoentes
    \item[$\square$] Verificar que cada grupo é dimensionless
    \item[$\square$] Validar resultado com físicaconhecida
    \item[$\square$] Documentar todas as suposições feitas
\end{itemize}

% =====================================================================
% APÊNDICE: EXERCÍCIOS
% =====================================================================
\appendix
\chapter{Exercícios Propostos}

\section{Nível 1: Básico}

\begin{problem}
Encontre a forma funcional da velocidade de uma partícula caindo em um fluido viscoso, que depende de:
- Massa: $m$
- Aceleração da gravidade: $g$
- Viscosidade: $\mu$
- Diâmetro da partícula: $d$
\end{problem}

\begin{problem}
Determine como a frequência de oscilação de um pêndulo de torsão depende de:
- Momento de inércia: $I$
- Rigidez torsional: $C$ (com dimensão $ML^2T^{-2}$)
\end{problem}

\begin{problem}
A amplitude de oscilação de um sistema amortecido diminui exponencialmente. Se a diminuição depende de:
- Coeficiente de amortecimento: $b$
- Massa: $m$
- Tempo: $t$

Encontre a forma da equação de decaimento.
\end{problem}

\section{Nível 2: Intermediário}

\begin{problem}
Usando o Teorema de Buckingham Pi, encontre os grupos adimensionais para um fluido em uma tubulação, com variáveis:
- Fluxo volumétrico: $Q$
- Diâmetro do tubo: $d$
- Queda de pressão por comprimento: $\nabla P$
- Viscosidade: $\mu$
- Densidade: $\rho$
\end{problem}

\begin{problem}
A resistência do ar em um objeto em queda depende de:
- Velocidade: $V$
- Área frontal: $A$
- Densidade do ar: $\rho$
- Viscosidade do ar: $\mu$
- Comprimento característico: $l$

Encontre todos os grupos dimensionless.
\end{problem}

\section{Nível 3: Avançado}

\begin{problem}
Para a transferência de calor por convecção, determine os grupos adimensionais envolvendo:
- Taxa de calor transferido: $\dot{Q}$
- Coeficiente de convecção: $h$
- Área: $A$
- Diferença de temperatura: $\Delta T$
- Condutividade térmica: $k$
- Viscosidade dinâmica: $\mu$
- Calor específico: $c_p$
\end{problem}

% =====================================================================
% REFERÊNCIAS
% =====================================================================
\begin{thebibliography}{99}

\bibitem{Dym2004} Dym, C. L. (2004). \textit{Principles of Mathematical Modeling}. 2nd Edition. Academic Press.

\bibitem{Buckingham1914} Buckingham, E. (1914). On Physically Similar Systems; Illustration of the Use of Dimensional Equations. \textit{Physical Review}, 4(4), 345--376.

\bibitem{Langhaar1951} Langhaar, H. L. (1951). \textit{Dimensional Analysis and Theory of Models}. John Wiley \& Sons.

\bibitem{Taylor1974} Taylor, E. S. (1974). \textit{Dimensional Analysis for Engineers}. Oxford University Press.

\bibitem{Barenblatt1996} Barenblatt, G. I. (1996). \textit{Scaling, Self-similarity, and Intermediate Asymptotics}. Cambridge University Press.

\end{thebibliography}

\end{document}
